% !TEX root = main.tex
\section{결론}
본 논문은 비공개 공급망 회계에서 사후 지정된 문제 입력의 후
속 사용을 제한하기 위한 zkDPP 구성을 제시하였다. 감사가 복
원하는 실제 관계와 원장의 소비 검사가 사용하는 식별자를 연
결하고, 분기·합류 뒤 도달하는 전체 미소비 객체를 대상으로 추

% Source PDF page 17, column 2
적·상태 변경· 완료 확인을 정의하였다. 이 구성은 기존의 비공
개 전이, 감사 암호화와 동결 기법을 생산·유통 객체의 생애주
기에 적용한 것이다.

조건부 증명은 정확성·건전성·프라이버시 세 성질로 구성하
였다. 정확성은 정상 조건에서 전이·공개 검증·추적·동결이 기
대 결과를 내는지를, 건전성은 실제 생성 값과 다른 전이나 잘
못된 추적 성공· 차단 완료를 받아들이게 할 수 있는지를 다룬
다. 두 키 암호화의 기밀성은 원장 프라이버시의 보조정리로 사
용하였다. 프라이버시 결과는 무작위 오라클 모형과 비밀·개방
값 공개가 없는 정직한 이력에 한정한다. 불변 소비 기록, 인증
된 시점 조회와 모든 소비 경로의 상태 검사가 각각 어떤 역할을
하는지 설명한다. 정확한 미소비 경계 집합 전체가 공통 시점에
미소비·차단 상태라면, 차단이 유지되는 동안 그 집합을 통한 추
가 소비는 거부된다. 반면 개별 동결의 성공 기록, 과거 시점의
정확한 추적 또는 일부 대상의 차단만으로 전체 실행의 완료가
따라오지는 않는다. 따라서 추적 정확성, 차단 상태와 완료 보
고를 구분해야 한다.

기호 모델에서는 네 사례군의 비어 있지 않은 각 접두 이력에
서 모든 기존 시작점에 대한 334개 추적 결과를 독립 그래프 계
산과 대조하였다. 이 모델의 record 시작 규칙은 규범의 seed
처리와 다르므로, 그 일치를 규범 알고리즘 전체의 적합성 증거
로 사용하지 않는다. 선택한 누락·부분 처리·경합 사례와 모든
소비 이벤트 종류의 차단 상태 검사도 재현하였다. 질량 0 출력
으로 탄소 반올림 잔여가 이동하는 허용 실행은 개별 회계식의
보존이 최종 제품군의 탄소 귀속을 자동 보장하지 않음을 보여
준다.

이 결과의 범위는 수락된 디지털 객체와 명시한 신뢰·원장
모델이다. 입력 증빙의 진실성, 이미 발급된 제품 주장의 취소
와 물리적 회수는 외부 문제로 남는다. 구체 암호 프로필의 보
안과 회로·계약·감사 클라이언트의 적합성은 별도로 검증해야
한다. 후속 평가에서는 그 검증과 함께 자료·키 가용성, 소비 경
합 및 실제 배포 비용을 확인할 필요가 있다.
